% Page 064 — page imprimée 61
% Suite de « MEYRUEIS DE 1942 à 1944 »

Et à partir de ce moment lundi 29 Mai, après l’inhumation des victimes de la Parade,
Meyrueis est \underline{surveillé}.

Le 6 Juin 1944 dès 5h du matin, 17 camions de 20 hommes en armes cernaient la ville. Mais
vers 9h00 ces troupes quittèrent rapidement Meyrueis. C’était le jour du débarquement. Sans
le débarquement que serait devenu Meyrueis ?

\begin{quote}
\textit{(en partant par la route de Florac ils s’arrêtèrent devant le moulin d’Ayres et nous les vîmes
descendre pour une pause « pipi », Damiani, Fafa et moi. Quelle peur rétrospective !)}
\end{quote}

\medskip

\noindent{\bfseries A JERUSALEM}

\medskip

Une avenue partant vers le désert est plantée d’arbres. Les Israéliens veulent inscrire sur les
troncs les noms de Français qui ont protégé les Juifs entre 40 et 45 en Europe.

\medskip

Il est indiqué par M. Frank Robert de Nimes ancien pasteur de Meyruies, de 41 à 45 des
noms de Meyrueisiens :

Les Dr. Bouisset et J. Cabanel, M. Rolland transporteur

Mme. Roussel du Château d’Ayres, château qui a été un centre de camouflage de première
grandeur. Propriétaire, avec sa fille Ginette Teissier du Cros et son mari Jean

M. Harrisson frère de Mme. Roussel, chef de réseau,. Des personnes en danger sont restées
plusieurs années au Château.

Des habitants des vallées et de Jontanels.

M. Cahen Salvador président de chambre au Conseil d’Etat.

M. Jacques Pfeiffer avocat à Paris.

\medskip

Notes prises à la suite d’une conversation entre Frank Robert et Marceline Causse de
Pradines. Non daté…

\medskip

\textit{Les notes en italique ont été rajoutées pour la compréhension du texte par Paul Loupiac\\
Le 5 Septembre 2002 au Moulin d’Ayres.}
